The anti-windup logic of PID loops was studied on the reverse-acting temperature loop of a three-state exothermic CSTR whose cooling valve is limited in amplitude and in slew rate, with six fixed-gain controllers as a test bed and true-temperature metrics. The contribution is an analysis of implementation, not a new anti-windup law. Proposition~\ref{prop:latch} holds for any reverse-acting loop around an open-loop-stable plant under its hypotheses, which the reactor is shown by simulation to reach; the regime map, the back-calculation constants and the loop-by-loop numbers refer to the reactor and valve studied. Four conclusions follow.
\begin{enumerate}
    \item \textbf{The anti-windup implementation can decide the verdict on a controller.} Sign-blind amplitude clamping, a frequent implementation error that reuses the direct-acting freeze condition, never freezes a reverse-acting integrator while it winds up and freezes it when it must unwind, which latches the valve at its upper stop (Proposition~\ref{prop:latch}). On the reactor the temperature settles $\vlatchgap$~K below the setpoint and the IAE of the nominal step rises from $\vZNdir$ to $\vZNsign$~K$\cdot$min for ZN and from $\vCCdir$ to $\vCCsign$~K$\cdot$min for CC. Comparisons that report the failure of a controller on a constrained loop should state their anti-windup logic (Section~\ref{subsec:checklist}).
    \item \textbf{The slew-rate trigger, not sign awareness, restores the loop, and its value depends on the valve.} Across three setpoint steps, a load profile, a valve deadband and a second operating point, the rate-aware trigger recovers the tracking and directional sign awareness is second order. Over slew limits of $20$ to $320$~L/min$^2$ the trigger matters for the empirical rules at every limit and for AMIGO and Optimal ITAE when the valve is slow, whereas it makes the slow, well-damped SIMC loop worse at every limit but the slowest; it is recommended when the loop demands faster moves than the valve allows.
    \item \textbf{Back-calculation with a well-chosen tracking constant matches the directional rule.} With one $T_t$ per loop inside $[T_d, T_i]$, chosen on the design step and judged on $\pm 5$~K, the worst IAE ratio to the directional rule lies between $\vttcvmin$ and $\vttcvmax$, and for the empirical rules back-calculation also moves the valve less and overshoots less in most cases; the usual defaults are up to $\vttsqrtmaxpct\%$ ($\sqrt{T_iT_d}$) and $\vtttimaxpct\%$ ($T_i$) worse. The directional rule needs no per-loop constant but depends on $S_{\max}$ and $\Delta t$.
    \item \textbf{The measurement path is part of the implementation.} The noise-to-valve gain of the empirical rules is set by $|K_p|(1+N)$ ($\vCSZN$ and $\vCSCC$~(L/min)/K for ZN and CC, against $\vCSSIMC$ for SIMC), not by a sensitivity peak, which on the exact linearization is only $\vSlinZN$ and $\vSlinCC$.
\end{enumerate}
The released code reproduces every table and figure and accepts new controllers behind the same interface. Future work will repeat the ablation on a second reverse-acting process, add valve stiction, and extend it to gain-scheduled controllers and to operating points inside the open-loop unstable region.
